% LaTeX preamble shared by the Design Booklet and every document PDF (passed to pandoc via -H).
\usepackage{fancyhdr}
\usepackage{emptypage}
\usepackage{amssymb}
\usepackage[dvipsnames]{xcolor}
\usepackage{ragged2e}
\usepackage{graphicx}
\usepackage[figuresright]{rotating}
\usepackage{xurl}

% ── Unicode symbols that appear in the FOC docs ──────────────────────────────
\usepackage{newunicodechar}
\newunicodechar{↔}{\ensuremath{\leftrightarrow}}
\newunicodechar{⇒}{\ensuremath{\Rightarrow}}
\newunicodechar{≈}{\ensuremath{\approx}}
\newunicodechar{✓}{\ensuremath{\checkmark}}
\newunicodechar{∈}{\ensuremath{\in}}
\newunicodechar{∇}{\ensuremath{\nabla}}
\newunicodechar{⋯}{\ensuremath{\cdots}}
\newunicodechar{ℝ}{\ensuremath{\mathbb{R}}}
\newunicodechar{ᵀ}{\ensuremath{^{\mathsf{T}}}}
\newunicodechar{β}{\ensuremath{\beta}}
\newunicodechar{θ}{\ensuremath{\theta}}
\newunicodechar{η}{\ensuremath{\eta}}
\newunicodechar{λ}{\ensuremath{\lambda}}
\newunicodechar{ε}{\ensuremath{\varepsilon}}
\newunicodechar{ω}{\ensuremath{\omega}}
\newunicodechar{ψ}{\ensuremath{\psi}}
\newunicodechar{α}{\ensuremath{\alpha}}
\newunicodechar{γ}{\ensuremath{\gamma}}
\newunicodechar{τ}{\ensuremath{\tau}}
\newunicodechar{σ}{\ensuremath{\sigma}}
\newunicodechar{φ}{\ensuremath{\varphi}}
\newunicodechar{Ω}{\ensuremath{\Omega}}
\newunicodechar{₀}{\textsubscript{0}}
\newunicodechar{₁}{\textsubscript{1}}
\newunicodechar{₂}{\textsubscript{2}}
\newunicodechar{ⁿ}{\textsuperscript{n}}
\newunicodechar{π}{\ensuremath{\pi}}
\newunicodechar{δ}{\ensuremath{\delta}}
\newunicodechar{ζ}{\ensuremath{\zeta}}
\newunicodechar{μ}{\ensuremath{\mu}}
\newunicodechar{ρ}{\ensuremath{\rho}}
\newunicodechar{≥}{\ensuremath{\geq}}
\newunicodechar{≤}{\ensuremath{\leq}}
\newunicodechar{≠}{\ensuremath{\neq}}
\newunicodechar{≈}{\ensuremath{\approx}}
\newunicodechar{⇄}{\ensuremath{\rightleftarrows}}
\newunicodechar{→}{\ensuremath{\rightarrow}}
\newunicodechar{←}{\ensuremath{\leftarrow}}
\newunicodechar{⋅}{\ensuremath{\cdot}}
\newunicodechar{×}{\ensuremath{\times}}
\newunicodechar{𝜃}{\ensuremath{\theta}}
\newunicodechar{𝜙}{\ensuremath{\phi}}
\newunicodechar{𝛷}{\ensuremath{\Phi}}
\newunicodechar{⁴}{\textsuperscript{4}}
\newunicodechar{⁵}{\textsuperscript{5}}
\newunicodechar{⁶}{\textsuperscript{6}}
\newunicodechar{⁷}{\textsuperscript{7}}
\newunicodechar{⁻}{\ensuremath{^{-}}}
\newunicodechar{⌈}{\ensuremath{\lceil}}
\newunicodechar{⌉}{\ensuremath{\rceil}}

% ── Bold math ───────────────────────────────────────────────────────────────
% Under unicode-math, \mathbf takes the bold text font, which has no Greek; \symbf takes the math font's.
\ifdefined\symbf
  \AtBeginDocument{\renewcommand{\mathbf}[1]{\symbf{#1}}}
\fi

% ── Running headers ──────────────────────────────────────────────────────────
\pagestyle{fancy}
\fancyhf{}
\renewcommand{\headrulewidth}{0.4pt}
\fancyhead[L]{\nouppercase{\leftmark}}
\fancyhead[R]{\thepage}

% ── Boxed equations (used for key results) ──────────────────────────────────
\usepackage{empheq}
\usepackage[most]{tcolorbox}
\tcbuselibrary{breakable}

% ── TikZ / pgfplots ─────────────────────────────────────────────────────────
\usepackage{tikz}
\usepackage{pgfplots}
\pgfplotsset{compat=1.18}
\usetikzlibrary{arrows.meta,positioning,calc}
\tikzset{arr/.style={-{Stealth[length=5pt]},thick}}
